\subsection{Finite support in the full metabelian quotient over number fields}

\begin{statement}{A.178}{Proved}
{A fixed abelian-cover quotient and the full metabelian quotient detect finite support}
Let \(k\) be a number field, let \(X/k\) be a smooth geometrically
connected hyperbolic curve, and let \(C\) be its smooth proper
compactification.  Put
\[
 \begin{gathered}
 E_X=\pi_1(X),\qquad \Pi=\pi_1(X_{\bar k}),\qquad
 \Pi'=\overline{[\Pi,\Pi]},\quad
 \Pi''=\overline{[\Pi',\Pi']},\\
 H_n=\ker\!\left(\Pi\longrightarrow
                         \Pi^{\rm ab}/n\Pi^{\rm ab}\right),
       \qquad H_n'=\overline{[H_n,H_n]}.
 \end{gathered}
\tag{A.178.1}
\]
All these subgroups used as quotient kernels are normal in the
indicated arithmetic group.  The following assertions concern
geometric conjugacy classes of continuous sections.

First, let \(Z\subset C\) be any finite subscheme.  Every section of
the actual arithmetic quotient \(E_X/H_3'\) which locally matches
actual point or cusp sections at every place, and whose chosen
proper local centres belong to \(Z(k_v)\) at a set of finite places
of Dirichlet density one, is the image of an actual full point or
cusp section over a point of \(Z(k)\).  That full point or cusp
class is unique.  In particular, the map from actual full point
and cusp classes to \(E_X/H_3'\) is injective.  This quotient
retains all coefficient primes in the open abelian group
\(H_3/H_3'\); it is not asserted to be a finite group.

Second, the same finite-support domain in the full geometrically
metabelian extension \(E_X/\Pi''\) consists exactly of the
images of the actual full point and cusp sections over \(Z(k)\).
Every finite local proper centre of such a section is the same
point of \(Z(k)\).  No rank or Tate--Shafarevich hypothesis is
imposed on \(C\), or on any covering curve.

For \(U_k=\mathbf P^1_k\setminus\{0,1,\infty\}\), this includes
the entire locally matching full-metabelian domain whose proper
local centres are cusps on a density-one set.  Every such section
is an actual global cusp image.  This conclusion does not require
its two global abelian coordinates to have been specified in
advance.

There is also an exact rational compact-Kummer fibre statement.
For \(g(C)\ge1\), choose \(R_0\in C(k)\) and put
\(j(R)=[R-R_0]\in J=\operatorname{Jac}(C)\).
For \(D\in J(k)\), the locally matching full-metabelian sections
with compact class \(\kappa(D)\) exist exactly for
\(D\in j(C(k))\).
Writing \(D=j(Q)\), their fibre is the single point section
when \(Q\in X(k)\), and the full cusp packet when
\(Q\in(C\setminus X)(k)\).
After choosing a rational parameter at this boundary point,
the latter fibre is bijective to \(\widehat{k^\times}\).
This assertion determines rational compact classes; no
classification of other completed Mordell--Weil classes is
included.
\end{statement}

\begin{proof}
\emph{The actual section covers and their proper genera.}
Write \(g=g(C)\), \(r=\#(C\setminus X)(\bar k)\).
For \(r>0\) the geometric abelianization is
\(\widehat{\mathbf Z}^{\,d}\), where \(d=2g+r-1\);
for \(r=0\), it is \(\widehat{\mathbf Z}^{\,2g}\).
These are the abelianizations of the actual surface presentations
\[
 \left\langle a_1,b_1,\ldots,a_g,b_g,c_1,\ldots,c_r\
 \middle|\
             \prod_{i=1}^g[a_i,b_i]\prod_{j=1}^r c_j=1
 \right\rangle^{\!\wedge},
\tag{A.178.2}
\]
where \(c_j\) is the meridian at the \(j\)-th geometric boundary
point; when \(r=0\), the boundary generators are absent.
The construction is invariant under the arithmetic outer action,
since \(H_n\) is characteristic in \(\Pi\).
Consequently \(H_n'\) is also arithmetic-normal, and
\(\Pi'\subset H_n\) gives the literal inclusions
\(\Pi''\subset H_n'\subset H_n\).

Let \(s_n:G_k\to E_X/H_n\) be the finite arithmetic projection of
the section under consideration.  Define, inside the original
arithmetic fundamental group,
\[
 E_n=\{e\in E_X:e\bmod H_n\in s_n(G_k)\}.
\tag{A.178.3}
\]
It is an open subgroup with
\(E_n\cap\Pi=H_n\) and \(E_n/H_n=G_k\).
The finite-étale-cover correspondence gives the actual cover
\(Y_n\to X\) with arithmetic group \(E_n\), and its geometric
cover has group \(\Pi/H_n\).  In particular \(Y_n\) is
geometrically connected and its constant field is \(k\).
Let \(C_n\) be the smooth proper normalization above \(C\).
The section in \(E_X/H_n'\) has its image in \(E_n/H_n'\);
this is its specified section on the geometrically abelianized
arithmetic group of \(Y_n\).

Here are the ramification and genus calculations, including the
one-puncture case.  Put \(D_n=[\Pi:H_n]\), and let \(r_n\) be
the geometric boundary size of \(Y_n\).  If \(r\ge2\), each
meridian in (A.178.2) is primitive in \(\Pi^{\rm ab}\), including
\(c_r=-\sum_{j<r}c_j\), so every boundary inertia has exact
order \(n\) in \(\Pi/H_n\).  If \(r=1\), its unique meridian
has abelian image zero.  Thus the actual ramification indices,
boundary counts, and Riemann--Hurwitz identities are
\[
 \begin{array}{c|c|c|c}
  &D_n&r_n&2g(C_n)-2\\ \hline
 r=0&n^{2g}&0&D_n(2g-2)\\
 r=1&n^{2g}&D_n&D_n(2g-2)\\
 r\ge2&n^{2g+r-1}&rD_n/n&
                 D_n(2g-2)+rD_n(1-1/n).
 \end{array}
\tag{A.178.4}
\]
In the last row, each of the \(rD_n/n\) boundary points has
ramification index \(n\); in the first two rows the proper cover
is unramified.  Hyperbolicity says \(2g-2+r>0\).  It follows
directly from (A.178.4) that for every multiple \(n\) of three
\[
       g(C_n)\ge1,\qquad
       r>0\ \Longrightarrow\ r_n\ge2.
\tag{A.178.5}
\]
For the smallest genus-zero case \(g=0,r=3,n=3\), the last
entry of (A.178.4) is zero, so \(g(C_3)=1\), and \(r_3=9\).
For \(r=1\), hyperbolicity gives \(g\ge1\), hence \(r_n=n^{2g}
\ge9\).  The proper case has \(g\ge2\).

We also specify the local fixed-sheet map.  A local geometric
conjugator in \(E_X/H_n'\) lifts to \(\Pi\).  Apply such a lift
to the actual full local point or cusp section whose quotient
is \(s\).  Its image modulo \(H_n\) is now literally \(s_n\),
so its full image lies in \(E_n\) after restriction to \(G_{k_v}\).
For a point section this is the fixed-sheet criterion on the
actual finite étale fibre, and gives a point of \(Y_n(k_v)\).
For a cusp section, the same inclusion in the boundary
decomposition group gives a \(G_{k_v}\)-fixed boundary point
of \(C_n\) and a full cusp section there: its residue action
is onto \(G_{k_v}\), so that boundary residue field is \(k_v\).
The inertia inclusion records the ramified tangent lift, not
just the proper point.  This argument also applies at real
places; the selected fixed sheet is retained.  Denote the
resulting proper point by \(P_{n,v}\in C_n(k_v)\).
Its image in \(C(k_v)\) is the originally selected centre.

\emph{Finite abelian descent before choosing any rational origin.}
Compactification gives a literal map
\[
 E_n/H_n'\longrightarrow
 \pi_1(C_n)\big/
          \overline{[\pi_1(C_{n,\bar k}),\pi_1(C_{n,\bar k})]}.
\tag{A.178.6}
\]
This map is obtained by killing boundary inertia and then taking
the geometric abelianization.  The section on the right has,
at every place, the point-section image of \(P_{n,v}\).

For clarity, the descent map used at this stage does not require
\(C_n(k)\) to be known nonempty.  Let \(F\) be a finite
commutative \(k\)-group scheme, and let
\(\xi\in H^1_{\acute et}(C_n,F)\) be a torsor.  In characteristic
zero it is finite étale.  Its fundamental-group cocycle has
coefficients \(F(\bar k)\) acted on through \(G_k\).
On the geometric group it is a homomorphism to the abelian
group \(F(\bar k)\), and therefore kills the closed geometric
commutator.  It descends through (A.178.6).  Pullback by the
global quotient section gives a literal class
\(\xi_s\in H^1(k,F)\), and the actual local section comparison
gives
\[
                       \xi(P_{n,v})=(\xi_s)_v
                                      \quad\hbox{for every }v.
\tag{A.178.7}
\]
Twisting the torsor by \(-\xi_s\) gives local points everywhere.
This is the implication (iii)\(\Rightarrow\)(i) of
Harari--Stix Theorem 2.1 with the closed geometric derived
subgroup.  The exact quotient is described in Remarks 2.2(3)
of the author version dated 11 February 2011, p.~8;
the same passage is the unnumbered Remarks (3) following
Theorem 2.1 in arXiv:1005.1302v2, dated 26 May 2010.
Both versions were checked
\cite[Theorem 2.1, Remarks 2.2(3), and Corollary 3.1
in the 2011 author version]{AgentHarariStix}.
The displayed cocycle map shows explicitly why it applies here.

Properness makes the family \(P_n=(P_{n,v})_v\) adelic.  At real
places use connected components, and omit the complex factors,
as in Stoll's modified adelic space.  Equation (A.178.7) proves
\[
                    P_n\in C_n(\mathbf A_k)^{\rm f\text{-}ab}_\bullet.
\tag{A.178.8}
\]
The preimage \(Z_n=(C_n\to C)^{-1}(Z)\) is finite, and the
assumed density-one support gives \(P_{n,v}\in Z_n(k_v)\) on
the same set of finite places.  Stoll's Theorem 8.2, with
\(g(C_n)\ge1\), now yields one actual point
\[
               Q_n\in Z_n(k),\qquad
               P_{n,v}=Q_n\quad\hbox{at every finite place}.
\tag{A.178.9}
\]
Its conclusion at real places is the modified-adelic component
statement.  We use only the displayed finite-place equalities
below.  The theorem does not assume a rational divisor class
of degree one, a rational origin, or a Tate--Shafarevich
finiteness statement
\cite[Theorem 8.2, p.~36 of the author version]{LegStoll2007}.
The conventions and density-one use of finite places agree
with the author's errata
\cite[Sections 1 and 3]{LegStollErrata2017}.

\emph{The exact affine-abelian recovery on one covering curve.}
We record all the remaining maps on a smooth affine or proper
curve \(Y/k\) with compactification \(\bar Y\), proper genus at
least one, and geometric boundary \(B\) either empty or of
size at least two.  Suppose its geometrically abelian section
has the local images of actual point or cusp sections, with
every finite proper centre equal to one \(Q\in\bar Y(k)\).
Write \(J=\operatorname{Jac}(\bar Y)\).  Abel--Jacobi with
origin \(Q\) is the actual morphism
\[
              j_Q:\bar Y\longrightarrow J,\qquad
                          R\longmapsto[R-Q].
\tag{A.178.10}
\]
It is injective on geometric points: a nonzero degree-one
divisor \(R-Q\) cannot be principal on a positive-genus curve,
as such a function would give a degree-one map to
\(\mathbf P^1\).

The point section at \(Q\) splits the proper arithmetic
abelian extension, whose geometric kernel is \(TJ\).
The difference of the proper section from that splitting
defines \(\eta\in H^1_{\rm cts}(k,TJ)\).
At each finite level \(m\), its local classes are Kummer
classes of the actual local points under (A.178.10);
thus the compatible finite classes give an element of
\(\varprojlim_m\operatorname{Sel}^{(m)}(k,J)\).
All its finite local values are zero, by the prescribed
proper centre.  Stoll's Proposition 3.7 gives zero in this
inverse Selmer group
\cite[Proposition 3.7, p.~10 of the author version]{LegStoll2007}.
This use has the precise inverse-Selmer domain; it asserts no
local-global injection on an arbitrary continuous cohomology
class.

The resulting continuous class is also zero.  Indeed, at each
level the set of zero-cochains which cobound the given finite
cocycle is a nonempty finite subset of \(J[m](\bar k)\).
The multiplication transition maps carry these sets into one
another.  Their inverse limit is nonempty by compactness,
and a compatible zero-cochain in \(TJ\) cobounds the full
continuous cocycle.  Thus the proper section is geometrically
conjugate to the point section at \(Q\).

The actual compactification homology sequence is
\[
 0\longrightarrow M_B\longrightarrow
             \pi_1(Y_{\bar k})^{\rm ab}
       \longrightarrow TJ\longrightarrow0,\qquad
 M_B=
 \frac{\widehat{\mathbf Z}(1)[B]}
      {\widehat{\mathbf Z}(1)\sum_{R\in B}e_R}.
\tag{A.178.11}
\]
For empty \(B\), put \(M_B=0\).  The maps are respectively
the boundary meridians and compactification; after splitting
the boundary, (A.178.2) proves that their sole boundary
relation is the displayed sum and that the quotient is the
compact homology.  The arithmetic action on \(M_B\) is the
actual boundary permutation times the cyclotomic character.
The compact homology is \(TJ\) by the Abel--Jacobi map
(A.178.10).  This is the same primitive boundary sequence
used in A.169.16, now with the specified \(k\) and \(\bar Y\).

The \(TJ\)-invariants vanish over \(k\) and over every finite
completion: local torsion of an abelian variety is finite, so
its inverse Tate system has zero invariant limit.  The global
assertion follows by restriction to one finite completion.
Consequently the cohomology map induced by \(M_B\) in
(A.178.11) is injective, both globally and at finite places.
After lifting the proper geometric conjugator, the difference
from the actual point section, or from a chosen full
tangential section at \(Q\), is therefore a uniquely determined
class
\[
                         \alpha\in H^1_{\rm cts}(k,M_B).
\tag{A.178.12}
\]

For \(Q\in Y(k)\), the local section is the actual point
section at \(Q\), so \(\alpha_v=0\) at every finite place.
Choose a finite Galois extension \(L/k\) splitting \(B\).
The restriction of \(M_B\) is a finite direct sum of
\(\widehat{\mathbf Z}(1)\).  The number-field completed
Kummer injection A.169.4 kills every coordinate of
\(\alpha|_L\).  Since \(M_B^{G_L}=0\), the explicit
restriction-injectivity argument following A.169.8 gives
\(\alpha=0\).  Thus the affine-abelian section is the
image of the actual full point section at \(Q\).

For \(Q\in B(k)\), its local sections are actual cusps at
that point.  Choose a parameter \(u_Q\in k(\bar Y)\) of
valuation one at \(Q\).  Its punctured formal disc gives
the actual inertia splitting, and its inertia line in
(A.178.11) is \(\widehat{\mathbf Z}(1)e_Q\).
Thus each \(\alpha_v\) belongs to the image of the
cohomology of that line.  The nonsplit boundary calculation
A.169.9--11 applies with the other boundary orbits unchanged:
over \(L\), all coordinates of \(\alpha\) in
\(\widehat{\mathbf Z}(1)[B\setminus\{Q\}]\) are the same
completed Kummer class; A.169.4 gives this equality globally
over \(L\), A.169.7 descends its scalar to \(k\), and
\(M_B^{G_L}=0\) gives equality before restriction.
It yields precisely one
\(\tau\in\widehat{k^\times}=H^1_{\rm cts}(k,\widehat{\mathbf Z}(1))\)
relative to the chosen tangent, and the actual full cusp
section with this parameter has the prescribed affine-abelian
image.

The uniqueness used here is literal.  In the cusp case,
\(e_Q\) is primitive because \(\#B\ge2\).  After splitting
\(B\), it is the negative diagonal in the free coordinates
indexed by \(B\setminus\{Q\}\).  The same Kummer and
restriction injections show that its map on global
\(H^1\) is injective.  Two full cusp splittings at that
point with the same parameter differ by an inertia
coboundary, hence are conjugate in its actual decomposition
group.  At distinct proper centres, the Abel--Jacobi
images are distinct by (A.178.10) and full local Kummer
injectivity at any finite place.  The latter follows from
the finite torsion and finite-rank \(\mathbf Z_p\)-module
structure of \(J(k_v)\): \(\bigcap_m mJ(k_v)=0\).
Therefore
\[
 \{\hbox{actual full point and cusp classes on }Y\}
       \longrightarrow
 \{\hbox{sections of }E_Y/\pi_1(Y_{\bar k})'\}
 \quad\hbox{is injective}.
\tag{A.178.13}
\]
This assertion concerns geometric conjugacy on that selected
curve.  It includes the disjointness of point and cusp
centres, and the injectivity of each full tangent parameter.

\emph{The single quotient \(E_X/H_3'\).}
Apply the preceding recovery to \(Y_3\), using
(A.178.5) and (A.178.9).  It supplies an actual full point
or cusp section \(\sigma_3\) of \(Y_3\) with the same
image as the given section in \(E_3/H_3'\).
The inclusion \(E_3\hookrightarrow E_X\) sends it to
an actual full point or cusp section on \(X\), centred at
the image of \(Q_3\in Z_3(k)\), hence in \(Z(k)\).
For a boundary point this pushforward is the literal
completed-parameter map
\[
 u_P=c\,u_{Q_3}^{\,e}v(u_{Q_3}),\qquad
 c\in k^\times,\quad v(0)=1,\qquad
                         \tau_P=c\,\tau_{Q_3}^{\,e}.
\tag{A.178.14}
\]
Here \(e\) is its actual ramification index.  The
principal formal factor has coherent roots of every order
in characteristic zero, so contributes zero to the
specialized Kummer class.  This proves the asserted
pushforward inside the actual cusp decomposition group.

We verify that this quotient separates full references on
\(X\), including its finite-sheet ambiguity.  Let two
actual full point or cusp sections have geometrically
conjugate images in \(E_X/H_3'\).  Lift that conjugator
to \(\Pi\) and conjugate one of the full sections.
Their finite arithmetic images in \(E_X/H_3\) are now
literally the same.  They consequently both lie in the
same actual arithmetic subgroup \(E_3\) of (A.178.3).
The fixed-sheet argument gives two actual full point or
cusp sections of this same \(Y_3\), and their images in
\(E_3/H_3'\) are equal.  Equation (A.178.13) makes them
conjugate by \(H_3\) as full sections on \(Y_3\).
After the inclusion into \(E_X\) and undoing the first
lifted conjugator, the original full sections are
geometrically conjugate on \(X\).  This proves the
injectivity and uniqueness in the first assertion; no
unproved identification of different deck sheets is used.

\emph{The full metabelian comparison with one fixed reference.}
Now start with a section of \(E_X/\Pi''\), and take the
cofinal sequence \(n=3j!\), \(j\ge1\).
Its arithmetic subgroups satisfy
\[
              n\mid m\quad\Longrightarrow\quad
                   E_m\subset E_n,\quad H_m'\subset H_n'.
\tag{A.178.15}
\]
These are actual cover maps \(Y_m\to Y_n\to X\);
their proper normalizations are finite and their boundary
maps are (A.178.14).  The preceding argument gives at
every level an actual full reference \(\sigma_n\) on
\(Y_n\), whose image in \(E_n/H_n'\) equals the given
section up to conjugacy by \(H_n/H_n'\).

Fix the full section \(\sigma=\sigma_3\) on \(Y_3\).
Push each \(\sigma_n\) through \(Y_n\to Y_3\).
Its affine-abelian image in \(E_3/H_3'\) is the same as
the given section and hence the same as that of \(\sigma\).
By the injectivity (A.178.13) on this fixed curve, the
pushforward is geometrically conjugate to that single
full \(\sigma\), by an element of \(H_3\).  This fixes
the full tangent, when the centre is a boundary point,
and not only the proper centre.
Combining this full conjugator with the preceding
\(H_n/H_n'\)-conjugator gives a conjugator in
\(H_3/H_n'\).  For every \(n\) define
\[
 \mathcal C_n=
 \left\{g\in H_3/\Pi'':
       g\,s(\gamma)\,g^{-1}=\sigma(\gamma)
       \text{ in }E_X/H_n'
                  \text{ for all }\gamma\in G_k\right\}.
\tag{A.178.16}
\]
Each set is nonempty by the surjection
\(H_3/\Pi''\to H_3/H_n'\), closed in a fixed compact
group, and the sets are nested under (A.178.15).

For completeness the exact group-theoretic kernel is
\[
                          \bigcap_{j\ge1}H_{3j!}'=\Pi''.
\tag{A.178.17}
\]
This proof works for the proper surface group as well as
for the free affine group.  If \(a\in\Pi'\setminus\Pi''\),
choose a continuous character \(f:\Pi'\to A\) to a finite
abelian group with \(f(a)\ne0\).  Choose an open normal
\(W\subset\Pi\) with \(W\cap\Pi'\subset\ker f\).
The formula
\[
                     \widetilde f(bw)=f(b),
                      \qquad b\in\Pi',\quad w\in W,
\tag{A.178.18}
\]
is well-defined and continuous on \(\Pi'W\).
It is multiplicative because every commutator
\([w,b]\) belongs to \(W\cap\Pi'\) and is killed by \(f\).
The open image of \(\Pi'W\) in the finitely generated
profinite abelian group \(\Pi^{\rm ab}\) contains
\(n\Pi^{\rm ab}\) for some \(n=3j!\).
Thus \(H_n\subset\Pi'W\) and \(\widetilde f\) kills
\(H_n'\), so \(a\notin H_n'\).
Elements outside \(\Pi'\) are outside every \(H_n'\).
The reverse inclusion is the already proved
\(\Pi''\subset H_n'\), giving (A.178.17).

Compactness now gives \(g\in\bigcap_n\mathcal C_n\).
Equation (A.178.17) makes its equalities actual equality
in \(E_X/\Pi''\).  This is the required single geometric
conjugator to the image of the full point or cusp section.
Its uniqueness as a full reference class follows from
the already proved injectivity into \(E_X/H_3'\).
Equation (A.178.9), pushed through the proper cover,
also gives equality of every finite local centre with
its one rational centre in \(Z(k)\).

\emph{The literal tripod specialization.}
For \(X=U_k\) the actual covers in (A.178.3) can be
written using the two finite Kummer coordinates of \(s\):
\[
 Y_n:\ A_n u^n+B_n v^n=1,\quad uv\ne0,\qquad
                      t=A_nu^n,\quad 1-t=B_nv^n.
\tag{A.178.19}
\]
Here \(A_n,B_n\in k^\times\) represent those classes.
The valuations at \(0\) and \(1\) show that the two
geometric Kummer characters are independent, so the
geometric degree is \(n^2\); the proper normalization
is the smooth plane curve
\(A_nU^n+B_nV^n=Z^n\).
For \(n\mid m\), the actual subgroup inclusion fixes
compatible roots and gives
\[
 (u_m,v_m)\longmapsto
     (c_{m,n}u_m^{m/n},d_{m,n}v_m^{m/n}),\qquad
 c_{m,n}^{\,n}=A_m/A_n,\quad
 d_{m,n}^{\,n}=B_m/B_n.
\tag{A.178.20}
\]
The constants lie in \(k\): the ratios of the corresponding
root functions have constant \(n\)-th powers, and their
geometrically connected function fields have constant
field \(k\).  Their compositions agree because these are
the selected inclusions \(E_m\subset E_n\), rather than
independently chosen sheets.
At \(n=3\), the open geometric abelian kernel has rank ten,
with its rank-eight boundary module and rank-two proper
Tate quotient.  Thus \(E_U/H_3'\) is an explicit arithmetic
quotient with geometric extension
\[
              1\longrightarrow\widehat{\mathbf Z}^{\,10}
                 \longrightarrow\Pi/H_3'
                 \longrightarrow(\mathbf Z/3)^2
                 \longrightarrow1.
\tag{A.178.21}
\]
All actions and extension data here are the actual ones
induced from \(E_U\); no splitting of this extension is
claimed.  Taking \(Z=\{0,1,\infty\}\) in the theorem
proves the stated whole cuspidal-local-centre consequence
over every number field.

\emph{The rational Abel--Jacobi fibres.}
Retain \(g(C)\ge1\), \(R_0\in C(k)\), and the actual morphism
\(j:C\to J\) in the statement.  In addition to the geometric
point injectivity proved at (A.178.10), this morphism is a
closed immersion.  Indeed \(h^0(C,\mathcal O_C(Q))=1\)
at every geometric point \(Q\), since any additional section
would give a degree-one function.  Riemann--Roch therefore
gives
\(h^0(C,\Omega_C^1(-Q))=g-1\).
The cotangent map of \(j\) is evaluation of the global
differentials at \(Q\), and is consequently surjective.
Thus \(j\) is unramified.  It is proper and geometrically
point-injective, hence finite onto its image with identical
residue fields; the surjection on cotangent spaces and
Nakayama's lemma make the local coordinate-ring maps
surjective.  This proves the asserted closed immersion.

The compact class of the section is supplied by the actual
maps
\[
 E_X/\Pi''\longrightarrow
 \pi_1(C)\big/\pi_1(C_{\bar k})'
       \xrightarrow[\ R_0\ ]{\ \sim\ }TJ\rtimes G_k,
 \qquad
                  [s]\longmapsto\eta_s\in H^1_{\rm cts}(k,TJ).
\tag{A.178.22}
\]
The first arrow is compactification and geometric
abelianization; the second is the splitting supplied by
\(R_0\) and its Abel--Jacobi map.  Let
\(\eta_s=\kappa(D)\) with \(D\in J(k)\).
For every finite place the local point/cusp comparison
and full local Kummer injection give
\(j(P_v)=D\) for its proper centre \(P_v\in C(k_v)\).
At even one such place the fibre
\[
                  W_D=\operatorname{Spec}k
                              \mathbin{\times}_{J,D}C
\tag{A.178.23}
\]
has a local point.  Since \(j\) is a closed immersion,
\(W_D\) is a closed subscheme of \(\operatorname{Spec}k\).
It is therefore empty or \(\operatorname{Spec}k\);
the local point excludes the empty case.
Thus \(D=j(Q)\) for a unique \(Q\in C(k)\).
Injectivity of \(j\) now gives \(P_v=Q\) at every
finite place.  The finite-support theorem with \(Z=\{Q\}\)
gives exactly the full point or cusp images stated.
Conversely each of those full references has compact
class \(\kappa(j(Q))\), by (A.178.22).

We finally check the full parameter assertion, including
possible finite-sheet conjugacy.  A rational parameter
\(u_Q\) splits the actual cusp group as
\(\widehat{\mathbf Z}(1)\rtimes G_k\); twisting that
splitting by the continuous Kummer cocycle of any
\(\tau\in\widehat{k^\times}\) constructs the full
cusp section.  Suppose two of them become conjugate
in the full metabelian quotient.  Project to
\(E_X/H_3'\) and use the lifted conjugator as in the
proof of (A.178.13) downstairs.  Both actual cusp
sections then lift to the same selected \(Y_3\).
Their equal affine-abelian images force the same
boundary point \(Q_3\) and the same full parameter
\(\tau_3\) there, by (A.178.13) on that curve.
The actual finite-cover parameter map (A.178.14), with
the same chosen \(u_Q\), gives for both sections
\[
                             \tau=c\,\tau_3^{\,e}.
\tag{A.178.24}
\]
Hence their original parameters are equal.
This proves the claimed bijection with
\(\widehat{k^\times}\) within the full metabelian fibre.
\end{proof}

The certificate
\path{code/metabelian_finite_support_certificate.py}
checks the actual Fermat-cover graph, its boundary meridian
lattice and proper homology rank, the displayed genus and
boundary formulas, and the polynomial identities for the
compatible tripod maps.  The finite computations supplement
the proof of the full inverse limit; they do not replace
Stoll's theorem or assert a finite-level replacement for the
all-coefficient kernel in (A.178.21).
